\input{preamble.tex}
\title{Gluing flag complexes along a tree}
\author{Human Compression project}
\date{Private draft, 11 October 2026}
\begin{document}
\maketitle

An abstract simplicial complex is \emph{flag} if each finite set of pairwise
adjacent vertices spans a simplex. A subcomplex is \emph{full} if it contains
every simplex of the ambient complex whose vertices it contains. We first
isolate a standard elementary gluing principle and prove the form needed
here. No novelty is claimed for the principle itself.

\begin{lemma}\label{lem:tree}
Let $T$ be a tree. For each vertex $v$ let $X_v$ be a flag simplicial
complex. For each edge $e=uv$ let $X_e$ be a complex equipped with
injective simplicial maps to $X_u$ and $X_v$, both with full image.
Form $X$ by identifying these images. Then each $X_v$ embeds as a full
subcomplex of $X$, and $X$ is flag.
\end{lemma}
\begin{proof}
An identification of vertices is a finite sequence of identifications
across edges of $T$. Whenever the underlying walk immediately crosses an
edge and returns, the two steps cancel, because the edge maps are
injective. Reducing the walk proves that a vertex of one $X_v$ cannot
be identified with a different vertex of the same $X_v$. It also proves
that the set $S(x)$ of vertices $v$ for which $x$ has a representative
in $X_v$ is a connected subtree: representatives along the unique path
between any two such vertices are forced and exist.

If two global vertices $x,y$ are represented in both $X_u$ and $X_v$,
they are represented in every block along the path from $u$ to $v$.
If they are adjacent in one of these blocks, fullness of the first edge
image puts their edge in that edge complex, and hence in the next block.
Repeating this argument transports adjacency along the path. More
generally, the identical argument transports every finite simplex.
Consequently every block is full in the union.

Now let $x_1,\ldots,x_m$ be pairwise adjacent vertices of $X$.
Their support subtrees $S(x_i)$ meet pairwise, since every edge of $X$
comes from a block. Finite families of pairwise intersecting subtrees
of a tree have a common vertex. Here is an explicit proof of that fact.
Root $T$ at any vertex. In each nonempty subtree take the unique vertex
nearest the root, and choose one, $p$, of greatest distance among these
finitely many nearest vertices. For another subtree choose a point where
it meets the subtree associated to $p$. The two nearest vertices lie on
the rooted path to that point; the deeper one is $p$. Connectedness of
the other subtree therefore puts $p$ in it. Thus $p$ lies in every subtree.

All $x_i$ are consequently represented in $X_p$. Adjacency transports
to that block as above. Since $X_p$ is flag, these representatives span
a simplex in $X_p$ and therefore in $X$.
\end{proof}

If $Y$ is flag and $\sigma$ is a simplex, its closed star is a full
subcomplex: its vertices are those of $\sigma$ and all vertices adjacent
to every vertex of $\sigma$. Any simplex on these vertices joins
$\sigma$ to a clique, and the flag property fills that clique.
This observation supplies the full edge inclusions in the next result.

\begin{theorem}\label{thm:amalgam}
Let $\alpha=(A,(A_i)_{i\in I_A})$ and
$\beta=(B,(B_i)_{i\in I_B})$ be flag-transitive coset geometries
of finite rank. Write $L=I_A\cap I_B$, $P=I_A\setminus L$ and
$Q=I_B\setminus L$. Suppose a fixed isomorphism identifies
$C=A_P$ with $B_Q$ and identifies $A_{P\cup J}$ with $B_{Q\cup J}$
for every $J\subseteq L$.
In $G=A*_C B$ put
\[
 G_i=\langle A_{\{i\}\cap I_A},B_{\{i\}\cap I_B}\rangle
 \quad(i\in I_A\cup I_B).
\]
Then the coset incidence system $(G,(G_i))$ is a flag-transitive
geometry. Its chamber stabilizer is the common subgroup
$H=A_{I_A}=B_{I_B}$, and for every $J\subseteq I_A\cup I_B$,
\[
 G_J=A_{J\cap I_A}*_{C\cap A_{J\cap I_A}}B_{J\cap I_B}.
\]
\end{theorem}
\begin{proof}
We use the usual tree of the amalgam: vertices are left cosets $gA,gB$,
edges are $gC$, and $gC$ joins $gA$ to $gB$. The reduced word theorem
proves that this graph is a tree. Using left cosets here only changes
the convention by inversion from the usual right-coset geometry.

At an $A$ vertex put the flag complex of $\alpha$ joined to a simplex
whose types are $Q$. At a $B$ vertex put the flag complex of $\beta$
joined to a simplex of types $P$. At an edge put the flag complex of
the common residue
\[
 (C,(C\cap A_l)_{l\in L})=(C,(C\cap B_l)_{l\in L})
\]
joined to the simplex of types $P\cup Q$. Flag transitivity identifies
this residue with the link of the standard $P$ face in $\alpha$ and
with the link of the standard $Q$ face in $\beta$. In each endpoint
block the edge complex is therefore the closed star of the corresponding
face, with the extra simplex already adjoined. These embeddings are full.
The construction is transported equivariantly to every coset vertex
and edge. Lemma~\ref{lem:tree} gives a flag complex $X$.

Every block chamber has one vertex of each type in $I_A\cup I_B$.
Flag transitivity in the two input geometries identifies its chambers
with the appropriate cosets of $H$. On an edge, the same identification
is by $C/H$, so these labels agree under every gluing. Thus the chambers
of $X$ are naturally $G/H$: equality of chamber labels is preserved
along the unique path of blocks supporting the chamber. In particular
$G$ acts transitively on them and their stabilizers are conjugates of $H$.

We give the same argument for an arbitrary standard face of type $J$,
including the empty face. Its stabilizer in the $A$ block is
$D_A=A_{J\cap I_A}$, its stabilizer in the $B$ block is
$D_B=B_{J\cap I_B}$, and its stabilizer on the edge is
$D_C=C\cap D_A=C\cap D_B$. The equality uses the fixed compatible
identification, not merely an abstract isomorphism of these groups.
The copies of this face that are identified across edges are exactly
the orbit of the standard copy under $\langle D_A,D_B\rangle$.
To justify this assertion locally, an edge $aC$ at the $A$ vertex
contains the standard $J$ face precisely when its $P$ face $aP$ is
incident with the standard $J\cap I_A$ face. Their union is then a
flag. Input flag transitivity supplies $d\in D_A$ carrying the
standard $P$ face to $aP$, so $aC=dC$. Thus the supporting incident
edges are exactly $D_A C/C$. At the other endpoint they are exactly
$D_B C/C$. A path transporting the face crosses only these edge
copies, so its consecutive changes of representatives are in $D_A$ or $D_B$.
Conversely each such change fixes the face. The reduced word theorem
identifies this generated subgroup with $D_A*_{D_C}D_B$.

Taking a one-vertex face identifies its stabilizer with the stated $G_i$.
Thus vertices of $X$ are precisely the corresponding cosets. Two vertices
in a simplex have a common chamber and hence a common coset representative;
conversely a common representative places them in a translate of the
standard chamber. Their incidence is therefore exactly coset incidence.
Taking the whole standard $J$ face gives the displayed intersection
formula, because its stabilizer is the intersection of its vertex
stabilizers. Finally every simplex of $X$ is contained in a block chamber,
so all flags extend to chambers. This proves the geometric assertion.
\end{proof}

\begin{corollary}\label{cor:residues}
Under the hypotheses of Theorem~\ref{thm:amalgam}, the amalgam geometry
is residually connected if and only if both input geometries are.
It is firm, thin or thick if and only if both inputs have the respective
property. Here the group actions need not be faithful.
\end{corollary}
\begin{proof}
The residue of a standard face of type $J$ has group $G_J$ and
parabolics $G_{J\cup\{i\}}$, for $i\notin J$. The intersection formula
therefore describes it by the same construction, using the corresponding
residues of the two input geometries. If $P\subseteq J$, its $A$ factor
is already in $C$ and the amalgam collapses to the $B$ factor. The
residue is exactly the residue of $\beta$ at $J\cap I_B$. If
$Q\subseteq J$, the symmetric statement holds. In particular, fixing
the standard $Q$ face recovers $\alpha$, and fixing the standard $P$
face recovers $\beta$; an empty such face simply means the whole geometry.
This proves necessity of residual connectivity.

For sufficiency, take a residue with at least two types left. The cases
just described reduce to residues of the inputs. Otherwise both
$P\setminus J$ and $Q\setminus J$ are nonempty. In the tree construction
of this residue, each $A$ block is a cone on its nonempty remaining
$Q$ simplex, and each $B$ block is a cone on its nonempty remaining
$P$ simplex. Each block is connected. The edge complexes contain the
nonempty simplex on $(P\cup Q)\setminus J$, so adjacent blocks meet.
Their union over a connected tree is connected. This includes all
residues of rank at least two and proves residual connectivity.

For a panel, $J=I\setminus\{i\}$ necessarily contains $P$ or $Q$.
The panel is consequently exactly an input panel by the first paragraph.
If $i\in L$, the two possible descriptions agree under the fixed
identification in $C$. Thus the panel sizes are the original indices
$[A_{I_A\setminus\{i\}}:H]$ or
$[B_{I_B\setminus\{i\}}:H]$, with every original index represented.
Comparison with $1$, $2$ or $3$ gives the three assertions.
\end{proof}

Theorem~\ref{thm:amalgam} repairs the forward flag-transitivity and
parabolic-intersection arguments for compatible parabolic amalgamations
in Piedade--Tranchida, arXiv:2505.24662v1, Section 3.2. The geometric
hypotheses are essential to this proof: arbitrary special subgroups
do not satisfy the corresponding intersection formula, as the separate
obstruction note shows. Corollary~\ref{cor:residues} supplies the resulting
residue and panel assertions. The converse flag-transitivity statement
for arbitrary input incidence systems and HNN extensions still require
their own arguments. This
draft has not yet had a distinct mathematical review.
\end{document}
